\unittitle{Shalom Aleichem\protect\footnote{These "angels of peace" descend from a story in the Talmud (Shabbat 119b): two angels, one good and one not, accompany each person home from synagogue on Shabbat eve. If the home is ready — candles lit, table set, peace between the people in it — the good angel blesses it, and even the other must answer "amen." The song is a welcome sung to unseen guests.}}{שָׁלוֹם עֲלֵיכֶם}
\kav{The time after shul, walking home in the darkness of the night, was one of fear and wonder, awe and giddiness. And so it became a custom for Jews to sing these words, asking that Holy Messengers protect us as we made our way home.}
\begin{hebtwocol}
\trow{Shalom aleichem, malachei hashalom}{שָׁלוֹם עֲלֵיכֶם, מַלְאֲכֵי הַשָּׁלוֹם}
\trow{Bo'achem l'shalom, malachei hashalom}{בּוֹאֲכֶם לְשָׁלוֹם, מַלְאֲכֵי הַשָּׁלוֹם}
\trow{Barchuni l'shalom, malachei ha'elyon}{בָּרְכוּנִי לְשָׁלוֹם, מַלְאֲכֵי הָעֶלְיוֹן}
\trow{B'tzetchem l'shalom, malachei hashalom}{בְּצֵאתְכֶם לְשָׁלוֹם, מַלְאֲכֵי הַשָּׁלוֹם}
\end{hebtwocol}
\vspace{4pt}
\translationrule
\transline{}{Peace upon you, angels of peace. Enter in peace, angels of peace. Bless me in peace, angels of the Most High. Go in peace, angels of peace.}
\unitsep

\unittitle{Yedid Nefesh\protect\footnote{Elazar Azikri wrote Yedid Nefesh in Safed in the 1500s and recorded it in his diary as a private mystical love poem, never meant for the prayerbook. It escaped his notebook and became one of the most beloved openings of Kabbalat Shabbat — a reminder that some of our most public liturgy began as someone's secret.}}{יְדִיד נֶפֶשׁ}
\kav{A 16th-century love song by R. Elazar Azikri of Safed — the soul chasing after Shabbat as the Beloved.}
\seclabel{V1}
\begin{hebtwocol}
\trow{Yedid nefesh, Av harachaman, meshoch avdecha el retzonecha.}{יְדִיד נֶפֶשׁ, אָב הָרַחֲמָן, מְשֹׁךְ עַבְדְּךָ אֶל רְצוֹנֶךָ.}
\trow{Yarutz avdecha kemo ayal, yishtachaveh el mul hadarecha.}{יָרוּץ עַבְדְּךָ כְּמוֹ אַיָּל, יִשְׁתַּחֲוֶה אֶל מוּל הֲדָרֶךָ.}
\trow{Ye'erav lo yedidutecha, minofet tzuf v'chol ta'am.}{יֶעֱרַב לוֹ יְדִידוּתֶךָ, מִנֹּפֶת צוּף וְכָל טָעַם.}
\end{hebtwocol}
\vspace{4pt}
\seclabel{V4}
\begin{hebtwocol}
\trow{Higaleh na ufros chaviv alai, et sukkat sh'lomecha.}{הִגָּלֶה נָא וּפְרֹשׂ חָבִיב עָלַי, אֶת סֻכַּת שְׁלוֹמֶךָ.}
\trow{Ta'ir eretz mik'vodecha, nagilah v'nism'cha vach.}{תָּאִיר אֶרֶץ מִכְּבוֹדֶךָ, נָגִילָה וְנִשְׂמְחָה בָךְ.}
\trow{Maher ahuv ki va mo'ed, v'chaneinu kimei olam.}{מַהֵר אֱהוֹב כִּי בָא מוֹעֵד, וְחָנֵּנוּ כִּימֵי עוֹלָם.}
\end{hebtwocol}
\vspace{4pt}
\translationrule
\transline{V1}{soul's beloved — compassionate one — pull me toward what you want. I'll run to you like a deer, fall at the light of your face. Your love is sweeter to me than honey, than anything.}
\transline{V4}{show yourself, please — spread over me, beloved, your canopy of peace. Let the whole earth shine with your presence; we'll dance and delight in you. Hurry, dear one — the time is here — be with us, like in the old days.}
\unitsep

\unittitle{Lechu Neranena}{לְכוּ נְרַנְּנָה}
\kav{The opening call of Kabbalat Shabbat — come, sing, shout to the rock of our liberation. We cross from week into Shabbat by raising our voices together. (Psalm 95)}
\begin{hebtwocol}
\trow{Lechu n'ran'na l'Adonai, nari'a l'tzur yish'einu.}{לְכוּ נְרַנְּנָה לַיהוָה, נָרִיעָה לְצוּר יִשְׁעֵנוּ.}
\trow{N'kad'ma fanav b'todah, biz'mirot nari'a lo.}{נְקַדְּמָה פָנָיו בְּתוֹדָה, בִּזְמִרוֹת נָרִיעַ לוֹ.}
\end{hebtwocol}
\vspace{4pt}
\translationrule
\transline{}{Come, let's sing out to HaShem — shout to the rock of our liberation. Let's come before HaShem's face with thanks, raise a cry with song.}
\unitsep

\unittitle{Shiru L'Adonai\protect\footnote{Twice in Kabbalat Shabbat we're commanded to sing a shir chadash — a "new song" (here and in Psalm 98). The rabbis noticed: why new? Because gratitude that only repeats itself calcifies. Every Shabbat asks for a song that has never been sung before, even to the same old melody.}}{שִׁירוּ לַיהוָה}
\kav{Singing is resistance work. Joy is resilience. Coming together is resilience. (Psalm 96)}
\begin{hebtwocol}
\trow{Shiru l'Adonai shir chadash, shiru l'Adonai kol ha'aretz.}{שִׁירוּ לַיהוָה שִׁיר חָדָשׁ, שִׁירוּ לַיהוָה כָּל־הָאָרֶץ.}
\trow{Shiru l'Adonai barchu shemo, bas'ru miyom l'yom y'shu'ato.}{שִׁירוּ לַיהוָה בָּרְכוּ שְׁמוֹ, בַּשְּׂרוּ מִיּוֹם־לְיוֹם יְשׁוּעָתוֹ.}
\trow{Sapru vagoyim k'vodo, b'chol ha'amim nifle'otav.}{סַפְּרוּ בַגּוֹיִם כְּבוֹדוֹ, בְּכָל־הָעַמִּים נִפְלְאוֹתָיו.}
\trow{Ki gadol Adonai um'hulal me'od, nora hu al kol elohim.}{כִּי גָדוֹל יְהוָה וּמְהֻלָּל מְאֹד, נוֹרָא הוּא עַל־כָּל־אֱלֹהִים.}
\trow{Ki kol elohei ha'amim elilim, v'Adonai shamayim asah.}{כִּי כָּל־אֱלֹהֵי הָעַמִּים אֱלִילִים, וַיהוָה שָׁמַיִם עָשָׂה.}
\trow{Hod v'hadar l'fanav, oz v'tif'eret b'mikdasho.}{הוֹד־וְהָדָר לְפָנָיו, עֹז וְתִפְאֶרֶת בְּמִקְדָּשׁוֹ.}
\trow{Havu l'Adonai mishp'chot amim, havu l'Adonai kavod va'oz.}{הָבוּ לַיהוָה מִשְׁפְּחוֹת עַמִּים, הָבוּ לַיהוָה כָּבוֹד וָעֹז.}
\trow{Havu l'Adonai k'vod sh'mo, s'u mincha u'vo'u l'chatzrotav.}{הָבוּ לַיהוָה כְּבוֹד שְׁמוֹ, שְׂאוּ־מִנְחָה וּבֹאוּ לְחַצְרוֹתָיו.}
\trow{Hishtachavu l'Adonai b'hadrat kodesh, chilu mipanav kol ha'aretz.}{הִשְׁתַּחֲווּ לַיהוָה בְּהַדְרַת־קֹדֶשׁ, חִילוּ מִפָּנָיו כָּל־הָאָרֶץ.}
\trow{Imru vagoyim Adonai malach, af tikon tevel bal timot, yadin amim b'meisharim.}{אִמְרוּ בַגּוֹיִם יְהוָה מָלָךְ, אַף־תִּכּוֹן תֵּבֵל בַּל־תִּמּוֹט, יָדִין עַמִּים בְּמֵישָׁרִים.}
\trow{Yis'm'chu hashamayim v'tagel ha'aretz, yir'am hayam um'lo'o.}{יִשְׂמְחוּ הַשָּׁמַיִם וְתָגֵל הָאָרֶץ, יִרְעַם הַיָּם וּמְלֹאוֹ.}
\trow{Ya'aloz sadai v'chol asher bo.}{יַעֲלֹז שָׂדַי וְכָל־אֲשֶׁר־בּוֹ.}
\end{hebtwocol}
\vspace{4pt}
\translationrule
\transline{V1}{Sing a new song — all the earth, sing. Sing, bless the name; carry the good news, day after day — liberation. Tell it among the nations: HaShem's glory; among all peoples — wonders. For HaShem is great, praised beyond measure — awesome above every power.}
\transline{V2}{All the gods of the nations are empty — but HaShem made the heavens. Splendor and glory before HaShem; strength and beauty in the holy place. Give to HaShem, you families of peoples — give HaShem glory and strength. Give HaShem the glory of the name; bring an offering, come into HaShem's courts.}
\transline{V3}{Bow before HaShem in holy splendor — tremble, all the earth. Tell the nations: HaShem reigns! The world stands firm, unshaken — peoples judged in fairness. The heavens rejoice, the earth is glad; the sea thunders, and every creature in it. The fields exult, and all that's in them.}
\unitsep

\unittitle{Mizmor Shiru L'Adonai}{מִזְמוֹר שִׁירוּ לַיהוָה}
\kav{Psalm 98 — a second new song. The whole world becomes a band: trumpets, shofar, rivers clapping, mountains singing together.}
\begin{hebtwocol}
\trow{Mizmor. Shiru l'Adonai shir chadash, ki nifla'ot asah,}{מִזְמוֹר, שִׁירוּ לַיהוָה שִׁיר חָדָשׁ, כִּי־נִפְלָאוֹת עָשָׂה,}
\trow{hoshi'a lo y'mino uzroa kodsho.}{הוֹשִׁיעָה־לּוֹ יְמִינוֹ וּזְרוֹעַ קָדְשׁוֹ.}
\trow{Hodia Adonai y'shu'ato, l'einei hagoyim gilah tzidkato.}{הוֹדִיעַ יְהוָה יְשׁוּעָתוֹ, לְעֵינֵי הַגּוֹיִם גִּלָּה צִדְקָתוֹ.}
\trow{Zachar chasdo ve'emunato l'veit Yisrael.}{זָכַר חַסְדּוֹ וֶאֱמוּנָתוֹ לְבֵית יִשְׂרָאֵל.}
\trow{Ra'u chol afsei aretz et y'shu'at Eloheinu.}{רָאוּ כָל־אַפְסֵי־אָרֶץ אֵת יְשׁוּעַת אֱלֹהֵינוּ.}
\trow{Hariu l'Adonai kol ha'aretz, pitzchu v'ran'nu v'zameiru.}{הָרִיעוּ לַיהוָה כָּל־הָאָרֶץ, פִּצְחוּ וְרַנְּנוּ וְזַמֵּרוּ.}
\trow{Zamru l'Adonai b'chinor, b'chinor v'kol zimrah.}{זַמְּרוּ לַיהוָה בְּכִנּוֹר, בְּכִנּוֹר וְקוֹל זִמְרָה.}
\trow{Ba'chatzotzrot v'kol shofar, hariu lifnei hamelech Adonai.}{בַּחֲצֹצְרוֹת וְקוֹל שׁוֹפָר, הָרִיעוּ לִפְנֵי הַמֶּלֶךְ יְהוָה.}
\trow{Yir'am hayam um'lo'o, tevel v'yoshvei vah.}{יִרְעַם הַיָּם וּמְלֹאוֹ, תֵּבֵל וְיֹשְׁבֵי בָהּ.}
\trow{N'harot yimcha'u chaf, yachad harim y'raneinu.}{נְהָרוֹת יִמְחֲאוּ־כָף, יַחַד הָרִים יְרַנֵּנוּ.}
\trow{Lifnei Adonai, ki va lishpot ha'aretz,}{לִפְנֵי־יְהוָה כִּי בָא לִשְׁפֹּט הָאָרֶץ,}
\trow{yishpot tevel b'tzedek, v'amim b'meisharim.}{יִשְׁפֹּט־תֵּבֵל בְּצֶדֶק, וְעַמִּים בְּמֵישָׁרִים.}
\end{hebtwocol}
\vspace{4pt}
\translationrule
\transline{V1}{A psalm. Sing HaShem a new song — wonders have been done. HaShem's right hand, holy arm — liberation. HaShem made the liberation known; before the eyes of the nations — justice revealed. HaShem remembered love and faithfulness to the house of Israel.}
\transline{V2}{All the ends of the earth have seen HaShem's liberation. All the earth — shout for HaShem — break open, sing, play on. Sing with the harp — harp and the sound of song. With trumpets, with the sound of the shofar — raise a cry before the sovereign.}
\transline{V3}{The sea thunders, and everything in it — the world and all who live in it. Rivers clap their hands; mountains sing together. Before HaShem, who comes to set the earth right — judging the world in justice, peoples in fairness.}
\unitsep

\unittitle{Mizmor L'David\protect\footnote{Psalm 29 names the divine voice — kol HaShem — seven times, a drumbeat that many read as the seven days of creation folded into a thunderstorm. Kabbalists paired it with the seven times we might otherwise say a name aloud, and gave it to Shabbat as the storm that clears the air before rest.}}{מִזְמוֹר לְדָוִד}
\kav{A psalm of the desert — a tribal cry from our nomadic ancestors gazing at the stars, their feet in the sand, howling at the storm. (Psalm 29)}
\begin{hebtwocol}
\trow{Mizmor l'David. Havu l'Adonai b'nei elim, havu l'Adonai kavod va'oz.}{מִזְמוֹר לְדָוִד. הָבוּ לַיהוָה בְּנֵי אֵלִים, הָבוּ לַיהוָה כָּבוֹד וָעֹז.}
\trow{Havu l'Adonai k'vod sh'mo, hishtachavu l'Adonai b'hadrat kodesh.}{הָבוּ לַיהוָה כְּבוֹד שְׁמוֹ, הִשְׁתַּחֲווּ לַיהוָה בְּהַדְרַת־קֹדֶשׁ.}
\trow{Kol Adonai al hamayim, el hakavod hir'im, Adonai al mayim rabim.}{קוֹל יְהוָה עַל־הַמָּיִם, אֵל־הַכָּבוֹד הִרְעִים, יְהוָה עַל־מַיִם רַבִּים.}
\trow{Kol Adonai ba'koach, kol Adonai behadar.}{קוֹל־יְהוָה בַּכֹּחַ, קוֹל יְהוָה בֶּהָדָר.}
\trow{Kol Adonai shover arazim, va'y'shaber Adonai et arzei ha'Levanon.}{קוֹל יְהוָה שֹׁבֵר אֲרָזִים, וַיְשַׁבֵּר יְהוָה אֶת־אַרְזֵי הַלְּבָנוֹן.}
\trow{Va'yarkidem k'mo egel, Levanon v'Siryon k'mo ven r'emim.}{וַיַּרְקִידֵם כְּמוֹ־עֵגֶל, לְבָנוֹן וְשִׂרְיֹן כְּמוֹ בֶן־רְאֵמִים.}
\trow{Kol Adonai chotzev lahavot eish.}{קוֹל־יְהוָה חֹצֵב לַהֲבוֹת אֵשׁ.}
\trow{Kol Adonai yachil midbar, yachil Adonai midbar Kadesh.}{קוֹל יְהוָה יָחִיל מִדְבָּר, יָחִיל יְהוָה מִדְבַּר קָדֵשׁ.}
\trow{Kol Adonai y'cholel ayalot va'yechesof y'arot, u'v'heichalo kulo omer kavod.}{קוֹל יְהוָה יְחוֹלֵל אַיָּלוֹת וַיֶּחֱשֹׂף יְעָרוֹת, וּבְהֵיכָלוֹ כֻּלּוֹ אֹמֵר כָּבוֹד.}
\trow{Adonai lamabul yashav, va'yeshev Adonai melech l'olam.}{יְהוָה לַמַּבּוּל יָשָׁב, וַיֵּשֶׁב יְהוָה מֶלֶךְ לְעוֹלָם.}
\trow{Adonai oz l'amo yiten, Adonai y'varech et amo vashalom.}{יְהוָה עֹז לְעַמּוֹ יִתֵּן, יְהוָה יְבָרֵךְ אֶת־עַמּוֹ בַשָּׁלוֹם.}
\end{hebtwocol}
\vspace{4pt}
\translationrule
\transline{V1}{A David song. You children of the gods — give to HaShem — give weight, give strength. Give the weight of the name; fall before HaShem in sacred trembling. HaShem's voice on the waters — the heavy one thunders — HaShem over the many waters.}
\transline{V2}{HaShem's voice — force. HaShem's voice — splendor. HaShem's voice shatters cedars — HaShem shatters the cedars of Lebanon. And makes them dance like a calf — Lebanon and Sirion leap like the young wild ox.}
\transline{V3}{HaShem's voice splits fire into flames. HaShem's voice — the wilderness writhes; HaShem shakes the wilderness of Kadesh. HaShem's voice sends the deer into labor, strips the forests bare — and in HaShem's shrine, all of it cries: kavod.}
\transline{Coda}{HaShem sat above the flood — and stays, present through all time. HaShem gives strength to the people — HaShem blesses the people with peace.}
\unitsep

\unittitle{L'cha Dodi\protect\footnote{L'cha Dodi was written by Shlomo HaLevi Alkabetz in 16th-century Safed, where Kabbalists would walk into the fields at dusk on Friday to greet Shabbat as a bride. The poem is an acrostic: the first letters of its verses spell out שלמה הלוי, the poet signing his own name inside the prayer.}}{לְכָה דוֹדִי}
\kav{A wedding poem by the mystics of Safed. They would dress in white and walk out to the fields to greet Shabbat like a bride, as the sun set over the mountains.}
\seclabel{Refrain}
\begin{hebtwocol}
\trow{Lecha dodi likrat kallah, p'nei Shabbat n'kab'la.}{לְכָה דוֹדִי לִקְרַאת כַּלָּה, פְּנֵי שַׁבָּת נְקַבְּלָה.}
\end{hebtwocol}
\vspace{4pt}
\seclabel{Shamor v'zachor}
\begin{hebtwocol}
\trow{Shamor v'zachor b'dibur echad, hishmi'anu El ham'yuchad.}{שָׁמוֹר וְזָכוֹר בְּדִבּוּר אֶחָד, הִשְׁמִיעָנוּ אֵל הַמְיֻחָד.}
\trow{Adonai echad u'shmo echad, l'shem u'l'tif'eret v'li't'hilah.    (L'cha…)}{יְהוָה אֶחָד וּשְׁמוֹ אֶחָד, לְשֵׁם וּלְתִפְאֶרֶת וְלִתְהִלָּה.    (לְכָה…)}
\end{hebtwocol}
\vspace{4pt}
\seclabel{Likrat Shabbat}
\begin{hebtwocol}
\trow{Likrat Shabbat l'chu v'nelcha, ki hi m'kor hab'racha.}{לִקְרַאת שַׁבָּת לְכוּ וְנֵלְכָה, כִּי הִיא מְקוֹר הַבְּרָכָה.}
\trow{Merosh mikedem n'sucha, sof ma'aseh b'machashava t'chila.    (L'cha…)}{מֵרֹאשׁ מִקֶּדֶם נְסוּכָה, סוֹף מַעֲשֶׂה בְּמַחֲשָׁבָה תְּחִלָּה.    (לְכָה…)}
\end{hebtwocol}
\vspace{4pt}
\seclabel{Hit'oreri}
\begin{hebtwocol}
\trow{Hit'oreri, hit'oreri, ki va orech kumi ori.}{הִתְעוֹרְרִי הִתְעוֹרְרִי, כִּי בָא אוֹרֵךְ קוּמִי אוֹרִי.}
\trow{Uri, uri, shir dabeiri, k'vod Adonai alayich nigla.    (L'cha…)}{עוּרִי עוּרִי שִׁיר דַּבֵּרִי, כְּבוֹד יְהוָה עָלַיִךְ נִגְלָה.    (לְכָה…)}
\end{hebtwocol}
\vspace{4pt}
\seclabel{Yamin usmol}
\begin{hebtwocol}
\trow{Yamin usmol tif'rotzi, v'et Adonai ta'aritzi.}{יָמִין וּשְׂמֹאל תִּפְרֹצִי, וְאֶת יְהוָה תַּעֲרִיצִי.}
\trow{Al yad ish ben Partzi, v'nismecha v'nagilah.    (L'cha…)}{עַל יַד אִישׁ בֶּן פַּרְצִי, וְנִשְׂמְחָה וְנָגִילָה.    (לְכָה…)}
\end{hebtwocol}
\vspace{4pt}
\seclabel{Bo'i v'shalom}
\begin{hebtwocol}
\trow{Bo'i v'shalom ateret ba'ala, gam b'simcha uv'tzahala.}{בּוֹאִי בְשָׁלוֹם עֲטֶרֶת בַּעְלָהּ, גַּם בְּשִׂמְחָה וּבְצָהֳלָה.}
\trow{Toch emunei am segula, bo'i chala, bo'i chala.    (L'cha…)}{תּוֹךְ אֱמוּנֵי עַם סְגֻלָּה, בּוֹאִי כַלָּה, בּוֹאִי כַלָּה.    (לְכָה…)}
\end{hebtwocol}
\vspace{4pt}
\translationrule
\transline{Refrain}{Come, beloved, to greet the bride — let's welcome the face of Shabbat.}
\transline{Shamor v'zachor}{"Protect" and "remember" — one word, one breath — the One let us hear it. HaShem is one, and the Name is one — for naming, for splendor, for praise.}
\transline{Likrat Shabbat}{Come, let's go toward Shabbat — she is the source of blessing. Poured out from time's beginning: last to be made, first to be imagined.}
\transline{Hit'oreri}{Wake up, wake up — your light has come — rise, shine. Awake, awake — break into song — HaShem's radiance breaks open on you.}
\transline{Yamin usmol}{Right and left you will burst through; it is HaShem you will adore. Through the one who breaks through — we will rejoice, we will be glad.}
\transline{Bo'i v'shalom}{Come in peace, crown of her beloved — with joy, with jubilation. Into the faithful, the treasured people — come, bride; come, bride.}
\footnote{The refrain's dodi — "my beloved" — is borrowed straight from the Song of Songs. Kabbalat Shabbat takes the Bible's most unabashed love poetry and points it at time itself: Shabbat as the beloved we've been waiting all week to meet.}
